% ---
% title: Concentration inequality
% date: 2025/12/19
% tag: ML
% ---
这是集中不等式，更准确地说，是 McDiarmid 不等式（有界差分不等式）的形式。

要看的式子是：

\[
\mathbb{P}\bigl[f(\mathcal X)-\mathbb{E}[f(\mathcal X)]\ge \epsilon\bigr] \le \exp\!\left( -\frac{2\epsilon^2}{\sum_{i=1}^m c_i^2} \right)
\]

它表达的意思是：

对于某个依赖于 \(m\) 个独立随机变量的函数 \(f(\mathcal X)\)，如果每个变量单独改变时对函数值的影响都不太大，那么 \(f(\mathcal X)\) 偏离其均值 \(\mathbb E[f(\mathcal X)] \) 超过 \(\epsilon\) 的概率会以指数速度衰减。

\hrule

\section{这是什么估计}

这是对尾概率

\[
\mathbb P\bigl(f(\mathcal X)-\mathbb E[f(\mathcal X)]\ge \epsilon\bigr)
\]

的一个上界估计。

也就是在估计：

\begin{quote}
“\(f(\mathcal X)\) 比它的平均值大至少 \(\epsilon\) 的概率有多大？”
\end{quote}

结论是它不超过

\[
\exp\!\left(-\frac{2\epsilon^2}{\sum_{i=1}^m c_i^2}\right).
\]

所以它是大偏差概率估计，也是集中性估计。

\hrule

\section{什么时候能用}

设\(\mathcal X=(X_1,\dots,X_m)\)

其中 \(X_1,\dots,X_m\) 相互独立，且函数 \(f\) 满足有界差分条件：

对于任意 \(i\)，如果只改变第 \(i\) 个变量，其他变量不变，那么函数值变化不超过 \(c_i\)，即

\[
|f(x_1,\dots,x_i,\dots,x_m)-f(x_1,\dots,x_i',\dots,x_m)|\le c_i.
\]

那么就有 \(McDiarmid\) 不等式。

这里的 \(c_i\) 就是“第 \(i\) 个变量对 \(f\) 的最大影响幅度”。

\hrule

\section{直观理解}

\(f\) 可能很复杂，但只要它对每个输入变量都“不太敏感”，并且这些输入变量彼此独立，\(f(\mathcal X)\) 就会高度集中在其期望附近。

也就是说，虽然 \(X_1,\dots,X_m\) 都是随机的，但 \(f(\mathcal X)\) 不会乱跑得太远。

\hrule

\section{这个上界怎么读}

上界是\(\exp\!\left(-\frac{2\epsilon^2}{\sum_{i=1}^m c_i^2}\right).\)

\(\epsilon\) 越大，概率越小；\(\sum c_i^2\) 越小，概率也越小。

也就是说，要求偏离均值很多时，这种事件很罕见；每个变量对结果的影响都很小时，整体就更稳定。

\hrule

\section{双边形式}

通常还会写成双边估计：

\[
\mathbb P\bigl(|f(\mathcal X)-\mathbb E[f(\mathcal X)]|\ge \epsilon\bigr) \le 2\exp\!\left(-\frac{2\epsilon^2}{\sum_{i=1}^m c_i^2}\right).
\]

开头给的是单边上尾版本。

\hrule

\section{一个简单例子}

如果\(f(X_1,\dots,X_m)=\sum_{i=1}^m X_i\)

并且每个 \(X_i\in[a_i,b_i]\)，那么改变一个 \(X_i\) 最多让和改变 \(b_i-a_i\)，于是可以取\(c_i=b_i-a_i.\)

这时 McDiarmid 不等式退化为与 Hoeffding 不等式很接近的形式。

\hrule

\section{在论文里怎么写}

在论文或证明里，一般可以写成：

\begin{quote}
由于 \(f(\mathcal X)\) 满足有界差分条件，且各随机变量独立，根据 McDiarmid 不等式，有
\end{quote}

\begin{quote}
\begin{quote}
\(\mathbb{P}\bigl[f(\mathcal X)-\mathbb E[f(\mathcal X)]\ge \epsilon\bigr] \le \exp\!\left(-\frac{2\epsilon^2}{\sum_{i=1}^m c_i^2}\right).\)
\end{quote}
\end{quote}

或者更口语一点：

\begin{quote}
该式给出了 \(f(\mathcal X)\) 偏离其期望的概率上界，表明当单个变量对函数值的影响受限时，\(f(\mathcal X)\) 会以指数型速度集中在其均值附近。
\end{quote}

\hrule

\section{实际怎么用}

通常是这几步：

\begin{enumerate}
  \item 写出随机变量 \(X_1,\dots,X_m\)；
  \item 定义目标函数 \(f(X_1,\dots,X_m)\)；
  \item 证明只改一个变量时，\(f\) 的变化不超过 c\_i；
  \item 代入 McDiarmid 不等式；
  \item 得到偏差概率上界。
\end{enumerate}
